% =============================================================================
% 82-lampiran-instrumen.tex — Lampiran C–F: Instrumen Psikometri
% =============================================================================

% ─────────────────────────────────────────────────────────────────────────────
\section*{Lampiran C: Instrumen PSS-10}

Perceived Stress Scale (PSS-10) dikembangkan oleh Cohen, Kamarck, dan Mermelstein (1983) dan terdiri dari sepuluh butir berskala Likert lima poin (0 = Tidak Pernah hingga 4 = Sangat Sering). Empat butir bersifat positif (butir 4, 5, 7, dan 8) dan diberi skor terbalik. Instrumen mengukur dua dimensi: ketidakberdayaan (butir 1, 2, 3, 6, 9, 10) dan efikasi diri (butir 4, 5, 7, 8). Skor total 0--40 diinterpretasikan menggunakan rentang yang umum digunakan: 0--13 (stres rendah), 14--26 (stres sedang), dan 27--40 (stres tinggi) (lih.\ Pieh dkk., 2021).

\noindent\textbf{Skala respons:} 0 = Tidak Pernah, 1 = Hampir Tidak Pernah, 2 = Kadang-kadang, 3 = Cukup Sering, 4 = Sangat Sering.

{\footnotesize
\renewcommand{\arraystretch}{1.3}
\begin{longtable}{|p{0.7cm}|p{8.8cm}|p{2.4cm}|p{1.4cm}|}
\hline
\textbf{No.} & \textbf{Pernyataan} & \textbf{Dimensi} & \textbf{Terbalik} \\
\hline
\endfirsthead
\hline
\textbf{No.} & \textbf{Pernyataan} & \textbf{Dimensi} & \textbf{Terbalik} \\
\hline
\endhead
\hline
\endfoot
1 & Seberapa sering kamu merasa kesal karena sesuatu yang terjadi secara tak terduga? & Helplessness & -- \\
\hline
2 & Seberapa sering kamu merasa tidak mampu mengendalikan hal-hal penting dalam hidupmu? & Helplessness & -- \\
\hline
3 & Seberapa sering kamu merasa gugup dan stres? & Helplessness & -- \\
\hline
4 & Seberapa sering kamu merasa yakin dengan kemampuanmu menangani masalah pribadi? & Self-Efficacy & Ya \\
\hline
5 & Seberapa sering kamu merasa segala sesuatu berjalan sesuai keinginanmu? & Self-Efficacy & Ya \\
\hline
6 & Seberapa sering kamu merasa tidak sanggup mengatasi semua hal yang harus kamu lakukan? & Helplessness & -- \\
\hline
7 & Seberapa sering kamu mampu mengendalikan rasa kesal dalam hidupmu? & Self-Efficacy & Ya \\
\hline
8 & Seberapa sering kamu merasa berada di atas kendali segalanya? & Self-Efficacy & Ya \\
\hline
9 & Seberapa sering kamu merasa marah karena hal-hal yang terjadi di luar kendalimu? & Helplessness & -- \\
\hline
10 & Seberapa sering kamu merasa kesulitan menumpuk begitu tinggi sehingga tidak bisa mengatasinya? & Helplessness & -- \\
\hline
\end{longtable}
}
\clearpage

% ─────────────────────────────────────────────────────────────────────────────
\section*{Lampiran D: Instrumen GPIUS-2}

Generalized Problematic Internet Use Scale 2 (GPIUS-2) dikembangkan oleh Caplan (2010). Versi yang digunakan merupakan adaptasi Bahasa Indonesia dengan skala Likert lima poin (1 = Sangat Tidak Sesuai hingga 5 = Sangat Sesuai). Lima belas butir terbagi ke dalam lima dimensi (POSI, regulasi suasana hati, preokupasi kognitif, penggunaan kompulsif, dan dampak negatif), masing-masing tiga butir, dengan dimensi orde-kedua disregulasi diri sebagai gabungan preokupasi kognitif dan penggunaan kompulsif. Karena GPIUS-2 bersifat dimensional dan tidak memiliki ambang klinis baku, rentang interpretasi rendah/sedang/tinggi yang ditampilkan platform bersifat operasional dan indikatif, mengacu pada nilai rata-rata normatif Indonesia (Reynaldo \& Sokang, 2016).

\noindent\textbf{Skala respons:} 1 = Sangat Tidak Sesuai, 2 = Tidak Sesuai, 3 = Netral, 4 = Sesuai, 5 = Sangat Sesuai.

{\footnotesize
\renewcommand{\arraystretch}{1.3}
\begin{longtable}{|p{0.7cm}|p{8.8cm}|p{2.4cm}|p{1.4cm}|}
\hline
\textbf{No.} & \textbf{Pernyataan} & \textbf{Dimensi} & \textbf{Terbalik} \\
\hline
\endfirsthead
\hline
\textbf{No.} & \textbf{Pernyataan} & \textbf{Dimensi} & \textbf{Terbalik} \\
\hline
\endhead
\hline
\endfoot
1 & Saya lebih memilih interaksi sosial secara daring daripada komunikasi tatap muka. & POSI & -- \\
\hline
2 & Saya telah menggunakan internet untuk berbicara dengan orang lain ketika merasa terisolasi. & MR & -- \\
\hline
3 & Saya merasa asyik mengingat pengalaman daring saya sebelumnya. & CP & -- \\
\hline
4 & Saya merasa sulit mengendalikan penggunaan internet saya. & CU & -- \\
\hline
5 & Penggunaan internet saya telah menyebabkan saya menemui kesulitan. & NO & -- \\
\hline
6 & Saya lebih aman berinteraksi dengan orang lain secara daring ketimbang secara langsung. & POSI & -- \\
\hline
7 & Saya telah menggunakan internet untuk merasa lebih baik ketika saya merasa sedih. & MR & -- \\
\hline
8 & Saya berpikir secara obsesif tentang online. & CP & -- \\
\hline
9 & Saya merasa bersalah atas waktu yang saya habiskan di internet. & CU & -- \\
\hline
10 & Penggunaan internet saya membuat saya sulit mengatur kehidupan saya. & NO & -- \\
\hline
11 & Saya lebih mudah untuk berkomunikasi dengan orang lain secara daring daripada secara langsung. & POSI & -- \\
\hline
12 & Saya telah menggunakan internet untuk membuat diri saya merasa lebih baik ketika saya merasa kecewa. & MR & -- \\
\hline
13 & Saya tidak bisa berhenti berpikir tentang menggunakan internet. & CP & -- \\
\hline
14 & Saya telah mencoba tanpa berhasil mengurangi berapa banyak waktu saya menggunakan internet. & CU & -- \\
\hline
15 & Saya telah melewatkan acara sosial atau kegiatan karena penggunaan internet saya. & NO & -- \\
\hline
\end{longtable}
}
\clearpage

% ─────────────────────────────────────────────────────────────────────────────
\section*{Lampiran E: Instrumen SRS}

Simplified Resilience Scale (SRS) dikembangkan oleh Manning, Carr, dan Kail (2016) sebagai ukuran ringkas resiliensi yang disusun dari butir Satisfaction With Life Scale (Diener dkk., 1985) dan Pearlin Mastery Scale (Pearlin \& Schooler, 1978), mengacu pada kerangka Resilience Scale (Wagnild \& Young, 1993). Versi asli terdiri dari dua belas butir dan bersifat unidimensional. Platform menggunakan adaptasi sebelas butir berskala Likert enam poin (1--6); lima butir (1, 3, 5, 7, 11) diberi skor terbalik. Pengelompokan ke dalam tiga aspek (efikasi, kepuasan, kontrol) bersifat deskriptif mengikuti instrumen sumber dan bukan struktur yang tervalidasi. Adaptasi ini belum melalui validasi psikometrik formal di Indonesia dan digunakan secara indikatif.

\noindent\textbf{Skala respons:} 1 = Sangat Tidak Setuju, 2 = Tidak Setuju, 3 = Agak Tidak Setuju, 4 = Agak Setuju, 5 = Setuju, 6 = Sangat Setuju.

{\footnotesize
\renewcommand{\arraystretch}{1.3}
\begin{longtable}{|p{0.7cm}|p{8.8cm}|p{2.4cm}|p{1.4cm}|}
\hline
\textbf{No.} & \textbf{Pernyataan} & \textbf{Dimensi} & \textbf{Terbalik} \\
\hline
\endfirsthead
\hline
\textbf{No.} & \textbf{Pernyataan} & \textbf{Dimensi} & \textbf{Terbalik} \\
\hline
\endhead
\hline
\endfoot
1 & Saya merasa tidak mungkin bagi saya mencapai tujuan yang saya perjuangkan. & Control & Ya \\
\hline
2 & Sejauh ini, saya sudah mendapatkan hal-hal penting yang saya inginkan dalam hidup. & Satisfaction & -- \\
\hline
3 & Jika sesuatu yang salah dapat terjadi pada saya, maka hal itu akan terjadi. & Control & Ya \\
\hline
4 & Saya puas dengan hidup saya. & Satisfaction & -- \\
\hline
5 & Apa yang terjadi dalam hidup saya sering kali berada di luar kendali saya. & Control & Ya \\
\hline
6 & Saya dapat melakukan hal-hal yang ingin saya lakukan. & Efficacy & -- \\
\hline
7 & Masa depan tampak tanpa harapan bagi saya dan saya tidak percaya bahwa segala sesuatu akan berubah menjadi lebih baik. & Control & Ya \\
\hline
8 & Ketika saya benar-benar ingin melakukan sesuatu, saya biasanya menemukan cara untuk berhasil melakukannya. & Efficacy & -- \\
\hline
9 & Dalam banyak hal, hidup saya hampir mendekati ideal. & Satisfaction & -- \\
\hline
10 & Saya bisa melakukan apa saja yang saya inginkan. & Efficacy & -- \\
\hline
11 & Tidak ada cara lain yang bisa saya lakukan untuk menyelesaikan masalah yang saya hadapi. & Control & Ya \\
\hline
\end{longtable}
}
\clearpage

% ─────────────────────────────────────────────────────────────────────────────
\section*{Lampiran F: Instrumen SRQ-29}

Self-Reporting Questionnaire (SRQ-29) dikembangkan oleh World Health Organization (Beusenberg \& Orley, 1994) dan terdiri dari dua puluh sembilan butir dengan respons biner (Ya/Tidak). Penilaian bersifat berbasis klaster: domain neurotik (butir 1--20) ditandai bila terdapat enam atau lebih jawaban `Ya'', sedangkan domain penggunaan zat (butir 21), psikotik (butir 22--24), dan stres pascatrauma (butir 25--29) ditandai bila terdapat minimal satu jawaban `Ya''. Instrumen ini bersifat skrining, bukan diagnostik; hasil positif memerlukan evaluasi klinis lanjutan. SRQ-29 telah divalidasi dan digunakan secara nasional di Indonesia (Riskesdas; Idaiani dkk., 2022).

\noindent\textbf{Skala respons:} 0 = Tidak, 1 = Ya.

{\footnotesize
\renewcommand{\arraystretch}{1.3}
\begin{longtable}{|p{0.7cm}|p{9.5cm}|p{3.0cm}|}
\hline
\textbf{No.} & \textbf{Pernyataan} & \textbf{Domain} \\
\hline
\endfirsthead
\hline
\textbf{No.} & \textbf{Pernyataan} & \textbf{Domain} \\
\hline
\endhead
\hline
\endfoot
1 & Apakah kamu sering merasa sakit kepala? & Neurotik \\
\hline
2 & Apakah kamu kehilangan nafsu makan? & Neurotik \\
\hline
3 & Apakah tidur kamu tidak nyenyak? & Neurotik \\
\hline
4 & Apakah kamu mudah merasa takut? & Neurotik \\
\hline
5 & Apakah kamu merasa cemas, tegang, atau khawatir? & Neurotik \\
\hline
6 & Apakah tangan kamu gemetar? & Neurotik \\
\hline
7 & Apakah kamu mengalami gangguan pencernaan? & Neurotik \\
\hline
8 & Apakah kamu merasa sulit berpikir jernih? & Neurotik \\
\hline
9 & Apakah kamu merasa tidak bahagia? & Neurotik \\
\hline
10 & Apakah kamu lebih sering menangis? & Neurotik \\
\hline
11 & Apakah kamu merasa sulit untuk menikmati aktivitas sehari-hari? & Neurotik \\
\hline
12 & Apakah kamu mengalami kesulitan untuk mengambil keputusan? & Neurotik \\
\hline
13 & Apakah aktivitas/tugas sehari-hari kamu terbengkalai? & Neurotik \\
\hline
14 & Apakah kamu merasa tidak mampu berperan dalam kehidupan ini? & Neurotik \\
\hline
15 & Apakah kamu kehilangan minat terhadap banyak hal? & Neurotik \\
\hline
16 & Apakah kamu merasa tidak berharga? & Neurotik \\
\hline
17 & Apakah kamu mempunyai pikiran untuk mengakhiri hidup kamu? & Neurotik \\
\hline
18 & Apakah kamu merasa lelah sepanjang waktu? & Neurotik \\
\hline
19 & Apakah kamu merasa tidak enak di perut? & Neurotik \\
\hline
20 & Apakah kamu mudah lelah? & Neurotik \\
\hline
21 & Apakah kamu minum alkohol lebih banyak dari biasanya atau apakah kamu menggunakan narkoba? & Penggunaan Zat \\
\hline
22 & Apakah kamu yakin bahwa seseorang mencoba mencelakai kamu dengan cara tertentu? & Psikotik \\
\hline
23 & Apakah ada yang mengganggu atau hal yang tidak biasa dalam pikiran kamu? & Psikotik \\
\hline
24 & Apakah kamu pernah mendengar suara tanpa tahu sumbernya atau yang orang lain tidak dapat mendengar? & Psikotik \\
\hline
25 & Apakah kamu mengalami mimpi yang mengganggu tentang suatu bencana/musibah atau adakah saat-saat kamu seolah mengalami kembali kejadian bencana itu? & PTSD \\
\hline
26 & Apakah kamu menghindari kegiatan, tempat, orang atau pikiran yang mengingatkan kamu akan bencana tersebut? & PTSD \\
\hline
27 & Apakah minat kamu terhadap teman dan kegiatan yang biasa kamu lakukan berkurang? & PTSD \\
\hline
28 & Apakah kamu merasa sangat terganggu jika berada dalam situasi yang mengingatkan kamu akan bencana atau jika kamu berpikir tentang bencana itu? & PTSD \\
\hline
29 & Apakah kamu kesulitan memahami atau mengekspresikan perasaan kamu? & PTSD \\
\hline
\end{longtable}
}
\clearpage
